%=====================================================================
\appendix
\section{The numerical companion}
\label{sec:numerics}

All figures are produced by \texttt{figures/spin\_boson\_numerics.py} (NumPy, SciPy, Matplotlib; about one minute on a laptop).
This appendix records what it computes, so that each curve can be traced to an equation in the text.

\paragraph{Model and parameters.}
$\omega_0=1$ (units), $S=\cos\theta\sz+\sin\theta\sx$ with $\theta=\pi/4$, Lorentzian bath with $\Omega=1$, $\gamma=0.5$, $T=0$, couplings $g=0.1$ and $0.3$; initial state $\ket{+x}$; time step $\dd t=0.02$, $t\le40$.
Kernels: \cref{eq:lorentz-kernels}.

\paragraph{Representation.}
Column-stacking vectorization, \cref{eq:vec}: $A^+=\Id\otimes A$, $A^-=A^T\otimes\Id$, $A^\cl=\tfrac12(A^++A^-)$, $A^\qu=A^+-A^-$, $\Lc=-iH^\qu$.
Bloch-basis matrix elements are $M_{\mu\nu}=\tfrac12\Tr(\sigma_\mu\mathcal M\sigma_\nu)$ (the coefficient convention of \cref{sec:lec1}).

\paragraph{Exact reference.}
The pseudomode master equation \eqref{eq:pseudomode} with $N=14$ oscillator levels (checked against $N=20$: maximal deviation $5\times10^{-15}$ in $\Pi(t)$ for $t\le4$ at $g=0.3$).
$\Pi(t)$ is obtained by propagating the four basis operators $\ket i\bra j\otimes\ket0\bra0$ with $\ee^{\Lc_{\rm tot}\dd t}$ and tracing out the oscillator.
The exact TCL generator is $\Gc(t)=\dot\Pi\Pi^{-1}-\Lc_S$ with $\dot\Pi$ computed from $\Lc_{\rm tot}$ (no numerical differentiation).
The exact NZ kernel is computed directly from its definition \eqref{eq:NZ}, $K(\tau)=\Pc\Lc_I\ee^{\Qc\Lc\tau}\Qc\Lc_I\Pc$, with $\Pc=\Tr_{\rm osc}(\cdot)\otimes\ket0\bra0$ built as a matrix on the $(2N)^2$-dimensional Liouville space of qubit plus oscillator; since the kernel is determined by $\Pi(t)$ alone (Exercise~\ref{ex:5-4}), it is also the kernel of the physical Lorentzian bath.

\paragraph{Approximations.}
$K^{(2)}(\tau)$ from \cref{eq:K2}; NZ2 solves the Volterra equation with the trapezoidal rule and an implicit treatment of $K(0)\rho_n$; TCL2 uses $\Gc^{(2)}(t)=\int_0^tK^{(2)}(\tau)\ee^{-\Lc_S\tau}\dd\tau$ (cumulative trapezoid) and the midpoint exponential; Markov uses $\Gc^{(2)}(t_{\max})$.

\paragraph{Checks printed by the script} (for $g=0.3$):
\begin{center}
\small
\begin{tabular}{@{}lll@{}}
\toprule
quantity & analytic (\cref{sec:lec5}) & numerical ($\Gc^{(2)}_\infty$)\\
\midrule
$\Gamma_\downarrow$, $\Gamma_\uparrow$ & $0.36000$, $0.00554$ & --- \\
$1/T_1=-\Gc_{zz}$ & $0.36554$ & $0.36552$\\
$1/T_2=-\tfrac12(\Gc_{xx}+\Gc_{yy})$ & $0.22512$ & $0.22512$\\
$\Gamma_\varphi=-\Gc_{xx}$ & $0.04235$ & $0.0424$\\
$\tilde\omega_0$ (secular / nonsecular) & $1.0222$ & $1.0219$\\
$r_z^{\rm eq}=\Gc_{z\Id}/(-\Gc_{zz})$ & $-0.96970$ & $-0.96970$\\
\bottomrule
\end{tabular}
\end{center}
The small differences are the $\dd t$ discretization and the nonsecular versus secular definitions.
The nonsecular generator has $\Gc_{xx}=-\Gamma_\varphi$ and $\Gc_{yy}=-\Gamma_\varphi-\Gamma_1$, whose average is $\Gamma_2$ (Example~5.3), and $\Gc_{yx}=2\delta\omega$, $\Gc_{xy}=0$.

\paragraph{Check of Theorem~\ref{thm:irreducible}.}
The function \texttt{check\_fourth\_order} evaluates the fourth-order irreducible kernel $K^{(4)}$ as a double time integral of $\Pc\Lc_I\ee^{\Lc_0\cdot}\Qc\Lc_I\ee^{\Lc_0\cdot}\Qc\Lc_I\ee^{\Lc_0\cdot}\Qc\Lc_I\Pc$ in the pseudomode space and compares $K-K^{(2)}$ with $K^{(4)}$ for $g=0.05,0.1,0.2$ (Example~5.1).

\paragraph{Check of Example~5.4.}
The function \texttt{check\_rwa\_kernel} builds the rotating-wave pseudomode model, $H_I=g(\sigma_+a+\sigma_-a^\dagger)$ with the same $\Omega=\omega_0$ and $\gamma$ (truncation $N=4$ suffices at $T=0$, where at most one excitation is ever present), evaluates $\Pc\Lc_I\ee^{\Qc\Lc\tau}\Qc\Lc_I\Pc$ literally, and compares the $ee$ and $eg$ elements with the closed forms \eqref{eq:Kee-exact-t} and \eqref{eq:Keg-exact} at $\tau\omega_0=0,0.5,1,2,4,8$ for $g=0.1,0.3$ (agreement $3\times10^{-17}$), then compares $K_{ee}-K^{(2)}_{ee}$ with the fourth-order term of \cref{eq:Kee-geometric} for $g=0.05,0.1,0.2$.

\paragraph{Check of \cref{sec:ill}.}
The function \texttt{check\_influence\_kernel} builds the discretized Feynman--Vernon path sum \eqref{eq:ill-U} in the $S$ eigenbasis for the Lorentzian bath ($\eta_n$ from \cref{eq:ill-eta}, $4^N$ paths) and prints: (i) $|U_N-\Pi(N\Delta t)|$ for $\Delta t=0.2,0.1,0.05$; (ii) the transfer tensors \eqref{eq:ill-ttm} computed from the exact $\Pi$ against the continuous kernel $K(n\Delta t)$; (iii) the transfer tensors of the Trotterized $U_N$ against the chain formula \eqref{eq:ill-chain} for $n\le3$; (iv) the inverse map \eqref{eq:ill-inverse}; (v) the Hamiltonian-learning numbers of Example~7.3, with Trotter-consistent and with exact input.
It runs in a few seconds.

\paragraph{Filter functions.}
\Cref{fig:filter} evaluates \cref{eq:filter} with the $1/\pi$ normalization (coupling $\tfrac12\sz\otimes B$) and $g=0.2$, for Ramsey, Hahn and CPMG-4 sequences.

\clearpage
